%! Semester 1, Units 1–3 — ACTUAL Final Exam (Answer Key)
% The scored cumulative final. Item for item it is the parallel form of ../practice_test/: the
% same three sections, the same 20 multiple-choice targets in the same order, the same five
% short-answer computations, and the same three free-response questions with the same parts and
% point values — every context different, so a student who worked the practice final has
% practised the moves and not memorised the answers. NEVER merged into any packet.
% TIME: written for 60–120 minutes (about 2 min per MC, 2 min per short answer, 12–15 min per
% free-response question: roughly 90 minutes for a typical student; 2 hours is the cap).
% WEIGHTS: Sections I and II follow the AP exam's unit weights (Unit 1 ~45%, Unit 3 ~35%,
% Unit 2 ~20%: 11 / 9 / 5 of the 25 items); Section III is one question per unit.
% POINTS: MC 20 x 2 = 40, short answer 5 x 3 = 15, free response 3 x 15 = 45; total 100.
%
% BLUEPRINT
%   MC 1 categorical vs quantitative (digits)      vet clinic records                    1.1
%   MC 2 groups of unequal size                    two high schools, part-time jobs      1.2
%   MC 3 IQR from a stem plot                      15 basketball games                   1.3, 1.5
%   MC 4 shape, and mean vs median                 60 urgent-care wait times             1.4, 1.5
%   MC 5 resistance                                shop salaries + the owner             1.5
%   MC 6 percentile                                a reading-test score                  1.5
%   MC 7 every section holds a quarter             80 days at a food truck               1.6
%   MC 8 parameter vs statistic                    city commute times                    1.7
%   MC 9 empirical rule                            bags of flour                         1.7
%   MC10 conditional relative frequency            plan x renewed                        2.1
%   MC11 what r says and does not say              distance from downtown vs rent        2.2
%   MC12 r from R-Sq, sign from the slope          used-car age vs price                 2.4
%   MC13 influential point                         employees vs electricity bill         2.5
%   MC14 which study is an experiment              five studies                          3.1
%   MC15 name the sampling method (cluster)        hotel chain                           3.2
%   MC16 name the bias and its direction           home internet, e-mail list            3.3
%   MC17 a larger sample does not fix bias         call-in vs SRS                        3.3
%   MC18 confounding                               optional tutoring                     3.4
%   MC19 randomized block design                   blood-pressure drug, two age groups   3.4
%   MC20 scope of inference (assigned, volunteers) stretching routine                    3.5
%   SA21 z-score                                   phone battery life                    1.7
%   SA22 upper outlier fence                       text messages sent                    1.6
%   SA23 residual                                  pizza delivery time                   2.3
%   SA24 count the treatments (two factors)        bean plants, water x soil             3.4
%   SA25 read a reshuffle dotplot                  100 reshuffles, observed 4            3.5
%   FRQ26 compare, outlier rule, resistance, empirical rule, z   coffee-shop service times  1.4–1.7
%   FRQ27 describe, equation + slope, residual, R-Sq, residual plot + extrapolation
%                                                  temperature vs iced coffees sold      2.2–2.5
%   FRQ28 SRS, nonresponse + direction, units/treatments/response, random assignment, scope
%                                                  a language-learning app               3.2–3.5
% ANSWERS (MC): 1C 2E 3A 4D 5B 6E 7C 8A 9D 10D 11B 12A 13E 14C 15B 16E 17D 18A 19C 20B
%
% THE DATA (every number checked)
%  MC2  North 120/400 = 0.30; South 100/250 = 0.40.
%  MC3  47 49 52 55 58 58 60 63 64 67 71 71 76 82 85: Q1 55, med 63, Q3 71, IQR 16, range 38.
%  MC4  bins of 10 min from 0: 18 20 11 6 3 1 1 (60), skewed right.
%  MC5  31 34 36 38 41: mean 36, median 36; with 120: mean 300/6 = 50, median 37.
%  MC7  10/18/22/34/60, n = 80, so 20 days per section.
%  MC9  N(5.00, 0.04): 4.92 = -2 SD, 5.04 = +1 SD -> 47.5% + 34% = 81.5%.
%  MC10 Basic 96/144 (240), Premium 120/40 (160), totals 216/184/400 -> 120/216 = 0.556.
%  MC12 n 25, xbar 6, s_x 2.5, ybar 14.5, s_y 4.2, r -0.80 -> b -1.344, a 22.56, S 2.574,
%       SE(b) 0.2102, T -6.39; SE(a) 1.362, T 16.57; R-Sq 64.0%, adj 62.4%.
%  MC13 x 2..12, y 15 14 19 18 23 22 27 25 31 30 34: b 1.94, r 0.970. With (24, 16): b 0.26, r 0.226.
%  SA21 (10.3 - 11.5)/0.8 = -1.5.   SA22 4/22/35/46/110: IQR 24, 46 + 36 = 82.
%  SA23 18.5 + 2.4(5) = 30.5; 27 - 30.5 = -3.5.   SA24 3 x 2 = 6.
%  SA25 100 reshuffles, step 1: -6:1 -5:1 -4:3 -3:7 -2:12 -1:17 0:20 1:16 2:11 3:7 4:3 5:1 6:1
%       (sum 100). At or above 4: 3 + 1 + 1 = 5.
%  FRQ26 drive-through 80/143/170/197/262 (IQR 54, fences 62 and 278: no outliers), about
%       N(170, 40). Walk-in 60/90/110/130, whisker to 185, outlier 240 (IQR 40, fence 190).
%       (d) 250 = 170 + 2(40) -> 2.5%. (e) (230-170)/40 = 1.5; (166-112)/27 = 2.0.
%  FRQ27 16 days, temperature x (F): 62 65 68 70 72 74 76 78 80 82 84 86 88 90 93 95
%        iced coffees y:             59 67 65 72 74 84 82 80 93 90 99 100 108 105 112 119
%        yhat = -51.98 + 1.774x, r 0.985, R-Sq 96.9%, adj 96.7%, S 3.267, SE(b) 0.08411,
%        T 21.09, SE(a) 6.690, T -7.77. Residuals: .99 3.66 -3.66 -.21 -1.75 4.70 -.85 -6.40
%        3.05 -3.50 1.96 -.59 3.86 -2.69 -1.01 2.44 (random scatter). (78, 80): yhat 86.39,
%        residual -6.39 with the printed coefficients. At 20 F: -16.5 coffees.
% NAMESTRIP: no name line — mirrors the blank (a graded page; the name is on the cover sheet only).
% NO teachernote here — scoring guidance belongs in the rubrics, never in a key.
% GENERATED FROM the blank.
% PAGINATION: each item carries a \boxguard sized to its height so no stem, figure or option
% list splits across a page; each section strip is guarded to travel with its first item.
% NO SKETCH ITEMS: every display is pre-drawn in TikZ and read.
% PAGE LOCKSTEP with ../../tests/actual_test/: every \writelines{n} here is answered by
% exactly n \ansline{} there, each terse enough not to wrap; every work block is byte-identical.
% Nothing else differs except the preamble, the header, and the \ans{} marks on the correct MC
% options.
\documentclass[12pt]{article}
\usepackage{apstats-article}
\usepackage{apstats-key}   % pulls in -boxes; do NOT also load -boxes
\usepackage{grading}
\gradingassignment{apstats-s01-actual-final}
\gradingcourse{advanced_statistics}

\begin{document}

\pageheader{Semester 1 \quad Units 1--3}{Final Exam --- Answer Key}
{\small\itshape Write each final answer in its box. Only what is inside a box is graded.}\par

\hfill\small Score:\;\blank{1.3cm}\;/\;100\normalsize
\vspace{0.06in}

\boxguard[8]
\begin{headlinebox}{sky}
\small
\textbf{You have up to 2 hours.} \textbf{Section I:} 20 multiple-choice questions, 40 points.
\textbf{Section II:} 5 short-answer questions, 15 points. \textbf{Section III:} 3 free-response
questions, 45 points. The exam covers Units 1, 2, and 3.
\end{headlinebox}

\vspace{0.08in}

\boxguard[10]
\begin{headlinebox}{goldbg}
\small
\textbf{Section I --- Multiple Choice} (2 points each). Write the letter of the single best answer in its Letter box. Read all
five options --- several of them name a real idea from the course, but only one fits \emph{this}
question.
\end{headlinebox}

\vspace{0.10in}

\begin{enumerate}[leftmargin=1.9em, itemsep=0.13in]

  % ── MC 1 ── 1.1
  \boxguard[8]
  \item A veterinary clinic records five variables for each patient: \emph{species},
        \emph{weight} (kilograms), the owner's \emph{ZIP code}, \emph{number of visits this
        year}, and \emph{age} (years).

        How many of these five variables are \textbf{quantitative}?\hfill\gradedright{\textit{Letter:}~\gradedblank[id=t1, type=choice, points=2, answer=C]{0.45in}{0.2in}}\rule[-0.2in]{0pt}{0.48in}
        \par\vspace{3pt}\noindent
        \hspace{1.2em}\textbf{(A)} 1 \hfill \textbf{(B)} 2 \hfill
        \textbf{(C)} \ans{3} \hfill \textbf{(D)} 4 \hfill \textbf{(E)} 5

  % ── MC 2 ── 1.2
  \boxguard[11]
  \item North High has 400 students, and 120 of them have a part-time job. South High has 250
        students, and 100 of them have a part-time job. Which statement is best supported?
        \begin{enumerate}[label=\textbf{(\Alph*)}, leftmargin=2.2em, itemsep=2pt]
          \item Jobs are more common at North, because 120 students is more than 100 students.
          \item Jobs are equally common, because each school has at least 100 students with one.
          \item No comparison is possible, because the schools are different sizes.
          \item Most South students have a part-time job.
          \item \ans{Jobs are more common at South: 40\% of its students have one, compared with}\\
                \ans{30\% of North's.}\hfill\gradedright{\textit{Letter:}~\gradedblank[id=t2, type=choice, points=2, answer=E]{0.45in}{0.2in}}\rule[-0.2in]{0pt}{0.48in}
        \end{enumerate}

  % ── MC 3 ── 1.3, 1.5
  \boxguard[13]
  \item The stem plot shows the points a basketball team scored in each of its 15 games.

        \vspace{3pt}
        {\centering\small
        \renewcommand{\arraystretch}{1.05}
        \begin{tabular}{@{} r | l @{}}
        4 & 7 \ 9 \\
        5 & 2 \ 5 \ 8 \ 8 \\
        6 & 0 \ 3 \ 4 \ 7 \\
        7 & 1 \ 1 \ 6 \\
        8 & 2 \ 5 \\
        \end{tabular}\qquad
        \raisebox{0.28in}{\footnotesize Key: $4 \mid 7$ means 47 points}\par}
        \vspace{2pt}

        What is the \textbf{interquartile range} of the points scored?\hfill\gradedright{\textit{Letter:}~\gradedblank[id=t3, type=choice, points=2, answer=A]{0.45in}{0.2in}}\rule[-0.2in]{0pt}{0.48in}
        \par\vspace{3pt}\noindent
        \hspace{1.2em}\textbf{(A)} \ans{16} \hfill \textbf{(B)} 19 \hfill
        \textbf{(C)} 24 \hfill \textbf{(D)} 38 \hfill \textbf{(E)} 63

  % ── MC 4 ── 1.4, 1.5
  \boxguard[18]
  \item The histogram shows how long each of 60 patients waited to be seen at an urgent-care
        clinic.

        \vspace{3pt}
        {\centering
        \begin{tikzpicture}[x=0.135cm, y=0.125cm]
          \foreach \y in {0,5,10,15,20} {
            \draw[linegray, very thin] (0,\y) -- (72,\y);
            \node[left, font=\tiny] at (0,\y) {\y};}
          \foreach \l/\c in {0/18, 10/20, 20/11, 30/6, 40/3, 50/1, 60/1}
            \draw[fill=skymid] (\l,0) rectangle (\l+10,\c);
          \draw[thick] (0,0) -- (72,0);
          \draw[thick] (0,0) -- (0,22);
          \foreach \x in {0,10,...,70} { \node[below, font=\tiny] at (\x,0) {\x}; }
          \node[font=\tiny] at (36,-5.0) {Wait time (minutes)};
          \node[rotate=90, font=\tiny] at (-6.5,11) {Patients};
        \end{tikzpicture}\par}
        \vspace{1pt}

        Which statement is most likely true?\hfill\gradedright{\textit{Letter:}~\gradedblank[id=t4, type=choice, points=2, answer=D]{0.45in}{0.2in}}\rule[-0.2in]{0pt}{0.48in}
        \begin{enumerate}[label=\textbf{(\Alph*)}, leftmargin=2.2em, itemsep=2pt]
          \item The distribution is skewed left, so the mean is likely less than the median.
          \item The distribution is skewed left, so the mean is likely greater than the median.
          \item The distribution is roughly symmetric, so the mean and median are about equal.
          \item \ans{The distribution is skewed right, so the mean is likely greater than the median.}
          \item The distribution is skewed right, so the mean is likely less than the median.
        \end{enumerate}

  % ── MC 5 ── 1.5
  \boxguard[11]
  \item The five employees of a small shop earn 31, 34, 36, 38, and 41 thousand dollars a
        year. The owner's salary of 120 thousand dollars is then added to the data set. What
        happens to the mean and the median?
        \begin{enumerate}[label=\textbf{(\Alph*)}, leftmargin=2.2em, itemsep=2pt]
          \item Both increase by the same amount.
          \item \ans{The mean increases from 36 to 50; the median increases only from 36 to 37.}
          \item The median increases from 36 to 50; the mean increases only from 36 to 37.
          \item Neither changes, because one value cannot move a summary of six values.
          \item The mean increases and the median decreases.\hfill\gradedright{\textit{Letter:}~\gradedblank[id=t5, type=choice, points=2, answer=B]{0.45in}{0.2in}}\rule[-0.2in]{0pt}{0.48in}
        \end{enumerate}

  % ── MC 6 ── 1.5
  \boxguard[10]
  \item Mia's score on a statewide reading test is at the 85th percentile. Which statement is
        correct?
        \begin{enumerate}[label=\textbf{(\Alph*)}, leftmargin=2.2em, itemsep=2pt]
          \item Mia answered 85\% of the questions correctly.
          \item Mia scored 85 points.
          \item 85\% of the students who took the test scored higher than Mia.
          \item Mia's score is 85\% of the highest score in the state.
          \item \ans{85\% of the students scored at or below Mia's score.}\hfill\gradedright{\textit{Letter:}~\gradedblank[id=t6, type=choice, points=2, answer=E]{0.45in}{0.2in}}\rule[-0.2in]{0pt}{0.48in}
        \end{enumerate}

  % ── MC 7 ── 1.6
  \boxguard[16]
  \item The boxplot shows the number of customers a food truck served on each of 80 days.

        \vspace{3pt}
        {\centering
        \begin{tikzpicture}[x=0.2cm, y=1cm]
          \draw[thick] (5,0) -- (65,0);
          \foreach \x in {10,20,...,60} {
            \draw (\x,-0.07) -- (\x,0.07); \node[below, font=\tiny] at (\x,-0.05) {\x};}
          \draw (10,0.45) -- (18,0.45);
          \draw (34,0.45) -- (60,0.45);
          \draw[fill=skymid] (18,0.25) rectangle (34,0.65);
          \draw[thick] (22,0.25) -- (22,0.65);
          \draw (10,0.33) -- (10,0.57);
          \draw (60,0.33) -- (60,0.57);
          \foreach \x in {10,18,22,34,60} \node[above, font=\tiny] at (\x,0.66) {\x};
          \node[font=\tiny] at (35,-0.62) {Customers per day};
        \end{tikzpicture}\par}
        \vspace{1pt}

        Which statement is true?\hfill\gradedright{\textit{Letter:}~\gradedblank[id=t7, type=choice, points=2, answer=C]{0.45in}{0.2in}}\rule[-0.2in]{0pt}{0.48in}
        \begin{enumerate}[label=\textbf{(\Alph*)}, leftmargin=2.2em, itemsep=2pt]
          \item More days had 34 to 60 customers than 18 to 22, because that whisker is longer.
          \item On about half of the days, the truck served fewer than 18 customers.
          \item \ans{About 20 days had 18 to 22 customers, and about 20 days had 34 to 60.}
          \item The mean number of customers per day was 22.
          \item The truck never served more than 34 customers in a day.
        \end{enumerate}

  % ── MC 8 ── 1.7
  \boxguard[11]
  \item A city wants to know the mean commute time of all 52{,}000 of its workers. It surveys
        a random sample of 400 workers and finds a mean commute of 27.3 minutes. Which
        statement is correct?
        \begin{enumerate}[label=\textbf{(\Alph*)}, leftmargin=2.2em, itemsep=2pt]
          \item \ans{27.3 minutes is a statistic; the mean commute of all 52{,}000 workers is}\\
                \ans{the parameter.}
          \item 27.3 minutes is a parameter, because it is a mean.
          \item 27.3 minutes is a parameter, because the sample was chosen at random.
          \item 400 is the statistic and 52{,}000 is the parameter.
          \item 27.3 minutes is both, because the sample is random.\hfill\gradedright{\textit{Letter:}~\gradedblank[id=t8, type=choice, points=2, answer=A]{0.45in}{0.2in}}\rule[-0.2in]{0pt}{0.48in}
        \end{enumerate}

  % ── MC 9 ── 1.7
  \boxguard[8]
  \item The weights of the bags of flour a machine fills are approximately normal, with mean
        5.00 pounds and standard deviation 0.04 pounds.

        About what percent of the bags weigh between 4.92 and 5.04 pounds?\hfill\gradedright{\textit{Letter:}~\gradedblank[id=t9, type=choice, points=2, answer=D]{0.45in}{0.2in}}\rule[-0.2in]{0pt}{0.48in}
        \par\vspace{3pt}\noindent
        \hspace{1.2em}\textbf{(A)} 34\% \hfill \textbf{(B)} 47.5\% \hfill
        \textbf{(C)} 68\% \hfill \textbf{(D)} \ans{81.5\%} \hfill \textbf{(E)} 95\%

  % ── MC 10 ── 2.1
  \boxguard[13]
  \item A streaming service classified 400 subscribers by their plan and by whether they
        renewed at the end of the year.

        \vspace{2pt}
        {\centering\small
        \renewcommand{\arraystretch}{1.15}
        \begin{tabular}{@{} l c c c @{}}
        \rowcolor{sky}
         & \textbf{Renewed} & \textbf{Did not renew} & \textbf{Total} \\
        \midrule
        Basic plan   & 96  & 144 & 240 \\
        Premium plan & 120 & 40  & 160 \\
        \midrule
        Total        & 216 & 184 & 400 \\
        \end{tabular}\par}
        \vspace{4pt}

        What proportion of the subscribers who \textbf{renewed} had the Premium plan?\hfill\gradedright{\textit{Letter:}~\gradedblank[id=t10, type=choice, points=2, answer=D]{0.45in}{0.2in}}\rule[-0.2in]{0pt}{0.48in}
        \par\vspace{3pt}\noindent
        \hspace{1.2em}\textbf{(A)} 0.30 \hfill \textbf{(B)} 0.40 \hfill
        \textbf{(C)} 0.54 \hfill \textbf{(D)} \ans{0.556} \hfill \textbf{(E)} 0.75

  % ── MC 11 ── 2.2
  \boxguard[11]
  \item For 30 apartments in one city, the correlation between distance from downtown (miles)
        and monthly rent (dollars) is $r = -0.82$. Which statement is correct?
        \begin{enumerate}[label=\textbf{(\Alph*)}, leftmargin=2.2em, itemsep=2pt]
          \item 82\% of the apartments that are farther from downtown have lower rent.
          \item \ans{There is a strong, negative linear association: apartments farther from
                downtown tend to have lower rent.}
          \item Moving an apartment farther from downtown causes its rent to fall.
          \item If distance were measured in kilometers, $r$ would change.
          \item Rent falls \$0.82 for each additional mile.\hfill\gradedright{\textit{Letter:}~\gradedblank[id=t11, type=choice, points=2, answer=B]{0.45in}{0.2in}}\rule[-0.2in]{0pt}{0.48in}
        \end{enumerate}

  % ── MC 12 ── 2.4
  \boxguard[11]
  \item A dealer regressed the price (thousands of dollars) of 25 used cars of one model on
        their age (years).

        \vspace{2pt}
        {\centering\footnotesize\ttfamily
        \begin{tabular}{@{} l r r r r @{}}
        Predictor & Coef & SE Coef & T & P \\ \hline
        Constant & 22.56 & 1.362 & 16.57 & 0.000 \\
        Age & -1.344 & 0.2102 & -6.39 & 0.000 \\
        \end{tabular}\par
        S = 2.574 \quad R-Sq = 64.0\% \quad R-Sq(adj) = 62.4\%\par}
        \vspace{2pt}

        What is the correlation between age and price?\hfill\gradedright{\textit{Letter:}~\gradedblank[id=t12, type=choice, points=2, answer=A]{0.45in}{0.2in}}\rule[-0.2in]{0pt}{0.48in}
        \par\vspace{3pt}\noindent
        \hspace{1.2em}\textbf{(A)} \ans{$-0.80$} \hfill \textbf{(B)} $-0.64$ \hfill
        \textbf{(C)} 0.41 \hfill \textbf{(D)} 0.64 \hfill \textbf{(E)} 0.80

  % ── MC 13 ── 2.5
  \boxguard[20]
  \item The scatterplot shows the number of employees and the monthly electricity bill
        (hundreds of dollars) for 12 small businesses. A least-squares line is fit to all 12.

        \vspace{3pt}
        {\centering
        \begin{tikzpicture}[xscale=0.30, yscale=0.105]
          \foreach \y in {10,20,30} {\draw[linegray, very thin] (0,\y) -- (26,\y);
            \node[left, font=\tiny] at (0,\y) {\y};}
          \draw[->, thick] (0,8) -- (27,8);
          \draw[->, thick] (0,8) -- (0,38);
          \foreach \x in {0,5,10,15,20,25} {\draw (\x,7.4) -- (\x,8.6);
            \node[below, font=\tiny] at (\x,7.6) {\x};}
          \foreach \x/\y in {2/15, 3/14, 4/19, 5/18, 6/23, 7/22, 8/27, 9/25, 10/31, 11/30,
                             12/34, 24/16}
            {\fill[navy] (\x,\y) circle[x radius=0.27, y radius=0.76];}
          \node[font=\tiny] at (13.5,1.2) {Employees};
          \node[rotate=90, font=\tiny] at (-3.4,23) {Bill (\$100s)};
        \end{tikzpicture}\par}
        \vspace{1pt}

        If the business with 24 employees were removed, what would happen?\hfill\gradedright{\textit{Letter:}~\gradedblank[id=t13, type=choice, points=2, answer=E]{0.45in}{0.2in}}\rule[-0.2in]{0pt}{0.48in}
        \begin{enumerate}[label=\textbf{(\Alph*)}, leftmargin=2.2em, itemsep=2pt]
          \item Nothing: one point out of 12 cannot change a least-squares line.
          \item The slope would decrease, and $r$ would move closer to 0.
          \item The slope would stay the same, and $r$ would become negative.
          \item The slope would decrease, and $r$ would move closer to 1.
          \item \ans{The slope would increase, and $r$ would move closer to 1.}
        \end{enumerate}

  % ── MC 14 ── 3.1
  \boxguard[14]
  \item Which of these studies is an \textbf{experiment}?
        \begin{enumerate}[label=\textbf{(\Alph*)}, leftmargin=2.2em, itemsep=2pt]
          \item A researcher surveys 200 students and compares the grades of those who eat
                breakfast with the grades of those who do not.
          \item A hospital uses its records to compare the recovery times of patients who chose
                one of two surgeries.
          \item \ans{A teacher flips a coin for each of 60 students to decide who studies with
                flashcards and who studies with practice tests, then compares quiz scores.}
          \item A sleep app tracks 5{,}000 users and compares coffee drinkers with non-drinkers.
          \item A pollster asks 1{,}000 randomly chosen voters about a new tax.\hfill\gradedright{\textit{Letter:}~\gradedblank[id=t14, type=choice, points=2, answer=C]{0.45in}{0.2in}}\rule[-0.2in]{0pt}{0.48in}
        \end{enumerate}

  % ── MC 15 ── 3.2
  \boxguard[13]
  \item A hotel chain has 60 hotels. To learn about guest satisfaction, it randomly selects 6
        of the hotels and surveys every guest staying in those 6 hotels tonight. What type of
        sample is this?
        \begin{enumerate}[label=\textbf{(\Alph*)}, leftmargin=2.2em, itemsep=2pt]
          \item A simple random sample, because the hotels were chosen at random.
          \item \ans{A cluster sample, because whole hotels were chosen at random and every guest
                in them was surveyed.}
          \item A stratified random sample, because the guests are divided into hotels.
          \item A convenience sample, because the guests were already at the hotels.
          \item A voluntary response sample, because guests can decline.\hfill\gradedright{\textit{Letter:}~\gradedblank[id=t15, type=choice, points=2, answer=B]{0.45in}{0.2in}}\rule[-0.2in]{0pt}{0.48in}
        \end{enumerate}

  % ── MC 16 ── 3.3
  \boxguard[13]
  \item To estimate the proportion of households in a county that have home internet service,
        a researcher e-mails a survey to a random sample of addresses from the county's
        e-mail newsletter list. Which is the main problem with this plan?
        \begin{enumerate}[label=\textbf{(\Alph*)}, leftmargin=2.2em, itemsep=2pt]
          \item Voluntary response bias, so the estimate is likely too low.
          \item Response bias, because the wording of the question is leading.
          \item There is no problem, because the sample was chosen at random.
          \item The sample is too small, so the estimate is likely too low.
          \item \ans{Undercoverage: households without internet service cannot be on the list,}\\
                \ans{so the estimate is likely too high.}\hfill\gradedright{\textit{Letter:}~\gradedblank[id=t16, type=choice, points=2, answer=E]{0.45in}{0.2in}}\rule[-0.2in]{0pt}{0.48in}
        \end{enumerate}

  % ── MC 17 ── 3.3
  \boxguard[14]
  \item A radio show asks listeners to call in with their opinion of a proposed toll road;
        12{,}000 call, and 74\% of them oppose it. A polling firm asks a simple random sample
        of 500 adults in the same city; 46\% oppose it. Which result is the more trustworthy
        estimate for the city's adults?
        \begin{enumerate}[label=\textbf{(\Alph*)}, leftmargin=2.2em, itemsep=2pt]
          \item The call-in result, because 12{,}000 is a much larger sample than 500.
          \item The call-in result, because callers care the most about the issue.
          \item Neither result, because neither one is a census.
          \item \ans{The SRS of 500: a larger sample does not remove the bias of a voluntary}\\
                \ans{response sample.}
          \item Both equally, because both come from the same city.\hfill\gradedright{\textit{Letter:}~\gradedblank[id=t17, type=choice, points=2, answer=D]{0.45in}{0.2in}}\rule[-0.2in]{0pt}{0.48in}
        \end{enumerate}

  % ── MC 18 ── 3.4
  \boxguard[13]
  \item A school offered optional after-school tutoring. Students who attended earned higher
        final grades, on average, than students who did not. The principal concludes that the
        tutoring raised grades. Which is the best criticism?
        \begin{enumerate}[label=\textbf{(\Alph*)}, leftmargin=2.2em, itemsep=2pt]
          \item \ans{Students chose whether to attend, so motivation is confounded with tutoring:
                motivated students may both attend and earn higher grades anyway.}
          \item The study had no placebo, so the grades cannot be compared.
          \item The students should have been chosen by a stratified random sample.
          \item The grades show response bias, because teachers assigned them.
          \item None is needed: the tutored students did better.\hfill\gradedright{\textit{Letter:}~\gradedblank[id=t18, type=choice, points=2, answer=A]{0.45in}{0.2in}}\rule[-0.2in]{0pt}{0.48in}
        \end{enumerate}

  % ── MC 19 ── 3.4
  \boxguard[13]
  \item A company tests a new blood-pressure drug against a placebo with 120 volunteers: 60 are
        under age 40, and 60 are age 40 or older. Age is known to affect blood pressure. Which
        plan is a \textbf{randomized block design}?
        \begin{enumerate}[label=\textbf{(\Alph*)}, leftmargin=2.2em, itemsep=2pt]
          \item Give the drug to everyone under 40 and the placebo to everyone 40 or older.
          \item Randomly assign 60 of the 120 volunteers to the drug and the rest to the placebo.
          \item \ans{Within each age group separately, randomly assign 30 volunteers to the drug and
                30 to the placebo.}
          \item Let each volunteer choose the drug or the placebo.
          \item Randomly select 60 volunteers and drop the rest.\hfill\gradedright{\textit{Letter:}~\gradedblank[id=t19, type=choice, points=2, answer=C]{0.45in}{0.2in}}\rule[-0.2in]{0pt}{0.48in}
        \end{enumerate}

  % ── MC 20 ── 3.5
  \boxguard[13]
  \item A researcher recruits 80 volunteers at one gym and randomly assigns 40 to a new
        stretching routine and 40 to their usual routine. After six months, the stretching
        group has significantly fewer injuries. Which conclusion is justified?
        \begin{enumerate}[label=\textbf{(\Alph*)}, leftmargin=2.2em, itemsep=2pt]
          \item The routine causes fewer injuries for all adults, because the result is
                significant.
          \item \ans{The routine causes fewer injuries for people like these volunteers, but the
                result cannot be generalized to all adults.}
          \item No cause can be claimed, because the subjects were volunteers.
          \item The result applies to all gym members, but no cause can be claimed.
          \item Nothing, because 80 people is too few.\hfill\gradedright{\textit{Letter:}~\gradedblank[id=t20, type=choice, points=2, answer=B]{0.45in}{0.2in}}\rule[-0.2in]{0pt}{0.48in}
        \end{enumerate}

\end{enumerate}

\vspace{0.14in}

\boxguard[22]
\begin{headlinebox}{goldbg}
\small
\textbf{Section II --- Short Answer} (3 points each). Work each problem in the space under it.
Write only the final number in its box. Only what is inside a box is graded.
\end{headlinebox}

\vspace{0.10in}

\begin{enumerate}[leftmargin=1.9em, itemsep=0.16in, resume]

  % ── SA 21 ── 1.7
  \boxguard[11]
  \item The battery life of one model of phone is approximately normal, with mean 11.5 hours
        and standard deviation 0.8 hours. One phone's battery lasted 10.3 hours. Find the
        $z$-score of this phone's battery life.\gradedbelow[r]{0.62in}{$z = $~\gradedblank[id=t21, type=short, points=3, answer=-1.5, check=numeric]{1.1in}{0.34in}}
        \begin{work}
          z &= \frac{10.3 - 11.5}{0.8} = \frac{-1.2}{0.8} = -1.5
        \end{work}
        \vspace{0.62in}

  % ── SA 22 ── 1.6
  \boxguard[11]
  \item The five-number summary of the number of text messages 40 students sent yesterday is
        4, 22, 35, 46, 110. By the $1.5 \times \text{IQR}$ rule, a value \emph{above} what
        number is an outlier?\gradedbelow[r]{0.62in}{Above~\gradedblank[id=t22, type=short, points=3, answer=82, check=numeric]{1.1in}{0.34in}}
        \begin{work}
          \text{IQR} &= 46 - 22 = 24 \\
          Q_3 + 1.5(\text{IQR}) &= 46 + 1.5(24) = 82
        \end{work}
        \vspace{0.50in}

  % ── SA 23 ── 2.3
  \boxguard[11]
  \item The line $\hat{y} = 18.5 + 2.4x$ predicts a pizza's delivery time $\hat{y}$ (minutes)
        from the distance $x$ (miles) to the customer. An order 5 miles away took 27 minutes.
        Find the residual for this order.\gradedbelow[r]{0.62in}{Residual $=$~\gradedblank[id=t23, type=short, points=3, answer=-3.5, check=numeric]{1.1in}{0.34in}}
        \begin{work}
          \hat{y} &= 18.5 + 2.4(5) = 30.5 \\
          \text{residual} &= y - \hat{y} = 27 - 30.5 = -3.5
        \end{work}
        \vspace{0.50in}

  % ── SA 24 ── 3.4
  \boxguard[10]
  \item An experiment on bean plants uses two factors: the amount of water (low, medium, or
        high) and the type of soil (sandy or clay). Every combination of water and soil is
        used. How many treatments are there?\gradedbelow[r]{0.50in}{Treatments:~\gradedblank[id=t24, type=short, points=3, answer=6, check=numeric]{1.1in}{0.34in}}
        \begin{work}
          3 \text{ amounts of water} \times 2 \text{ types of soil} &= 6 \text{ treatments}
        \end{work}
        \vspace{0.60in}

  % ── SA 25 ── 3.5
  \boxguard[20]
  \item In an experiment, the mean score of the treatment group was \textbf{4 points higher}
        than the mean score of the control group. A computer reshuffled the scores into two new
        groups 100 times, as if the treatment made no difference. The dotplot shows the 100
        simulated differences. How many of the 100 reshuffles gave a difference of 4 points or
        more?
        \begin{center}
        \begin{tikzpicture}[x=0.82cm, y=0.115cm]
          \draw[thick] (-6.6,0) -- (6.6,0);
          \foreach \x in {-6,...,6} {
            \draw (\x,0) -- (\x,-1.3); \node[below, font=\tiny] at (\x,-1.1) {$\x$};}
          \foreach \x/\n in {-6/1,-5/1,-4/3,-3/7,-2/12,-1/17,0/20,1/16,2/11,3/7,4/3,5/1,6/1}
            \foreach \k in {1,...,\n} \fill[navy] (\x,\k) circle (1.3pt);
          \node[font=\tiny] at (0,-6.6) {Simulated difference in mean score (points), treatment $-$ control};
          \node[font=\scriptsize\bfseries] at (0,24.5) {100 random reshuffles};
        \end{tikzpicture}
        \end{center}
        \vspace{-6pt}
        \begin{work}
          3 + 1 + 1 &= 5
        \end{work}
        \hfill Reshuffles at 4 or more:~\gradedblank[id=t25, type=short, points=3, answer=5, check=numeric]{1.1in}{0.34in}

\end{enumerate}

\vspace{0.14in}

\boxguard[34]
\begin{headlinebox}{goldbg}
\small
\textbf{Section III --- Free Response} (15 points each). Each part has its own box: write your
final answer to the part inside it. Show calculations in the space above the box. Answer
\emph{in context} --- name the people or things studied, the variables, and their units.
\end{headlinebox}

\vspace{0.10in}

\begin{enumerate}[leftmargin=1.9em, itemsep=0.16in, resume]

  % ── FRQ 26 ── Unit 1
  \boxguard[32]
  \item \textbf{(15 points)} The owner of a coffee shop recorded the service time (seconds)
        for 40 drive-through customers and for 40 walk-in customers.

        \vspace{4pt}
        \noindent
        \begin{minipage}[c]{0.55\linewidth}\centering
        \begin{tikzpicture}[x=0.036cm, y=1.15cm]
          \draw[thick] (40,0) -- (280,0);
          \foreach \x in {50,100,...,250} {
            \draw (\x,-0.07) -- (\x,0.07); \node[below, font=\tiny] at (\x,-0.05) {\x};}
          % drive-through 80/143/170/197/262
          \draw (80,1.15) -- (143,1.15); \draw (197,1.15) -- (262,1.15);
          \draw[fill=skymid] (143,0.95) rectangle (197,1.35); \draw[thick] (170,0.95) -- (170,1.35);
          \draw (80,1.03) -- (80,1.27); \draw (262,1.03) -- (262,1.27);
          \node[font=\tiny, anchor=south west] at (40,1.38) {Drive-through};
          % walk-in 60/90/110/130, whisker 185, outlier 240
          \draw (60,0.45) -- (90,0.45); \draw (130,0.45) -- (185,0.45);
          \draw[fill=goldbg] (90,0.25) rectangle (130,0.65); \draw[thick] (110,0.25) -- (110,0.65);
          \draw (60,0.33) -- (60,0.57); \draw (185,0.33) -- (185,0.57);
          \fill[navy] (240,0.45) circle (1.6pt);
          \node[font=\tiny, anchor=south west] at (40,0.68) {Walk-in};
          \node[font=\tiny] at (160,-0.55) {Service time (seconds)};
        \end{tikzpicture}
        \end{minipage}\hfill
        \begin{minipage}[c]{0.42\linewidth}\centering\footnotesize
        \renewcommand{\arraystretch}{1.15}
        \setlength{\tabcolsep}{3.5pt}
        \begin{tabular}{@{} l c c c c c @{}}
        \rowcolor{sky}
         & \textbf{Min} & $\boldsymbol{Q_1}$ & \textbf{Med} & $\boldsymbol{Q_3}$ & \textbf{Max} \\
        \midrule
        Drive-through & 80 & 143 & 170 & 197 & 262 \\
        Walk-in       & 60 & 90  & 110 & 130 & 240 \\
        \end{tabular}
        \end{minipage}

        \vspace{6pt}
        \textbf{(a)} Compare the distributions of service time for drive-through and walk-in
        customers.\\[4pt]
        \gradedhere{\gradedblank[id=t26a, type=rubric, points=3, rubric=../rubrics/apstats-s01-actual-final-t26a.yaml]{\linewidth+6pt}{\gradedlinesheight{3}}}\ansline{Drive-through times are typically longer: median 170 s vs.\ 110 s for walk-in.}\\
        \ansline{Drive-through times also vary more: IQR 54 s vs.\ 40 s for walk-in.}\\
        \ansline{Walk-in times have one high outlier (240 s); drive-through times have none.}\\
        \ansline{}\\
        \ansline{}\\
        \ansline{}\\
        \ansline{}\\
        \ansline{}

        \vspace{6pt}
        \boxguard[20]
        \textbf{(b)} Use the $1.5 \times \text{IQR}$ rule to show that the walk-in time of 240
        seconds is an outlier.
        \begin{work}
          \text{IQR} &= 130 - 90 = 40 \text{ seconds} \\
          Q_3 + 1.5(\text{IQR}) &= 130 + 1.5(40) = 190 \text{ seconds}
        \end{work}
        \strut\gradedhere{\gradedblank[id=t26b, type=rubric, points=3, rubric=../rubrics/apstats-s01-actual-final-t26b.yaml]{\linewidth+6pt}{\gradedlinesheight{3}}}\ansline{240 seconds is greater than the upper fence of 190 seconds, so the walk-in time}\\
        \ansline{of 240 seconds is an outlier.}\\
        \ansline{}\\
        \ansline{}\\
        \ansline{}\\
        \ansline{}\\
        \ansline{}\\
        \ansline{}

        \vspace{6pt}
        \boxguard[16]
        \textbf{(c)} The owner finds that the 240 was a recording error: that walk-in customer
        really waited 140 seconds. Without calculating, explain how correcting it changes the
        mean and the median of the walk-in times.\\[4pt]
        \gradedhere{\gradedblank[id=t26c, type=rubric, points=3, rubric=../rubrics/apstats-s01-actual-final-t26c.yaml]{\linewidth+6pt}{\gradedlinesheight{3}}}\ansline{The mean decreases (by $100/40 = 2.5$ s): it uses every value, so it is not resistant.}\\
        \ansline{The median stays 110 s: 140 s is still above the median, so the middle of the}\\
        \ansline{ordered walk-in times does not move. The median is resistant.}\\
        \ansline{}\\
        \ansline{}\\
        \ansline{}\\
        \ansline{}\\
        \ansline{}

        \vspace{6pt}
        \boxguard[16]
        \textbf{(d)} Drive-through service times are approximately normal, with mean 170 seconds
        and standard deviation 40 seconds. About what percent of drive-through customers are
        served in \emph{more than} 250 seconds? Explain.\\[4pt]
        \gradedhere{\gradedblank[id=t26d, type=rubric, points=3, rubric=../rubrics/apstats-s01-actual-final-t26d.yaml]{\linewidth+6pt}{\gradedlinesheight{3}}}\ansline{250 s $= 170 + 2(40)$, two standard deviations above the mean. About 95\% of}\\
        \ansline{drive-through times are within 2 SDs of the mean, leaving 5\% outside, half on}\\
        \ansline{each side: about 2.5\% of drive-through customers are served in more than 250 s.}\\
        \ansline{}\\
        \ansline{}\\
        \ansline{}\\
        \ansline{}\\
        \ansline{}

        \vspace{6pt}
        \boxguard[22]
        \textbf{(e)} After the correction, walk-in times have mean 112 seconds and standard
        deviation 27 seconds. Which is more unusual \emph{for its own kind of customer}: a
        drive-through time of 230 seconds or a walk-in time of 166 seconds? Justify with
        $z$-scores.
        \begin{work}
          z_{\text{drive-through}} &= \frac{230 - 170}{40} = 1.5 \\[3pt]
          z_{\text{walk-in}} &= \frac{166 - 112}{27} = 2.0
        \end{work}
        \strut\gradedhere{\gradedblank[id=t26e, type=rubric, points=3, rubric=../rubrics/apstats-s01-actual-final-t26e.yaml]{\linewidth+6pt}{\gradedlinesheight{3}}}\ansline{The walk-in time of 166 s is more unusual: it is 2.0 standard deviations above}\\
        \ansline{the walk-in mean, but 230 s is only 1.5 SDs above the drive-through mean.}\\
        \ansline{}\\
        \ansline{}\\
        \ansline{}\\
        \ansline{}\\
        \ansline{}\\
        \ansline{}

  % ── FRQ 27 ── Unit 2
  \boxguard[36]
  \item \textbf{(15 points)} A bakery owner recorded the daily high temperature (degrees
        Fahrenheit) and the number of iced coffees sold on 16 summer days. The scatterplot,
        computer output, and residual plot are below.

        \vspace{4pt}
        \noindent
        \begin{minipage}[b]{0.49\linewidth}\centering
        \begin{tikzpicture}[xscale=0.165, yscale=0.045]
          \foreach \y in {60,80,100,120} {\draw[linegray, very thin] (58,\y) -- (98,\y);
            \node[left, font=\tiny] at (58,\y) {\y};}
          \draw[->, thick] (58,50) -- (99,50);
          \draw[->, thick] (58,50) -- (58,127);
          \foreach \x in {60,70,80,90} {\draw (\x,48.8) -- (\x,51.2);
            \node[below, font=\tiny] at (\x,49.4) {\x};}
          \foreach \x/\y in {62/59, 65/67, 68/65, 70/72, 72/74, 74/84, 76/82, 78/80, 80/93,
                             82/90, 84/99, 86/100, 88/108, 90/105, 93/112, 95/119}
            {\fill[navy] (\x,\y) circle[x radius=0.39, y radius=1.43];}
          \node[font=\tiny] at (78,36) {Temperature ($^\circ$F)};
          \node[rotate=90, font=\tiny] at (51.5,88) {Iced coffees sold};
        \end{tikzpicture}
        \end{minipage}\hfill
        \begin{minipage}[b]{0.49\linewidth}\centering
        \begin{tikzpicture}[xscale=0.165, yscale=0.215]
          \draw[->, thick] (58,-8) -- (58,8);
          \draw[thick] (58,-8) -- (99,-8);
          \draw[dashed, slate] (58,0) -- (99,0);
          \foreach \y in {-6,-3,0,3,6} {\draw (57.6,\y) -- (58.4,\y);
            \node[left, font=\tiny] at (58,\y) {\y};}
          \foreach \x in {60,70,80,90} {\draw (\x,-8.25) -- (\x,-7.75);
            \node[below, font=\tiny] at (\x,-8.1) {\x};}
          \foreach \x/\e in {62/0.99, 65/3.66, 68/-3.66, 70/-0.21, 72/-1.75, 74/4.70, 76/-0.85,
                             78/-6.40, 80/3.05, 82/-3.50, 84/1.96, 86/-0.59, 88/3.86, 90/-2.69,
                             93/-1.01, 95/2.44}
            {\fill[navy] (\x,\e) circle[x radius=0.39, y radius=0.30];}
          \node[font=\tiny] at (78,-11) {Temperature ($^\circ$F)};
          \node[rotate=90, font=\tiny] at (51.8,0) {Residual (coffees)};
        \end{tikzpicture}
        \end{minipage}

        \vspace{4pt}
        {\centering\footnotesize\ttfamily
        \begin{tabular}{@{} l r r r r @{}}
        Predictor & Coef & SE Coef & T & P \\ \hline
        Constant & -51.98 & 6.690 & -7.77 & 0.000 \\
        Temperature & 1.774 & 0.08411 & 21.09 & 0.000 \\
        \end{tabular}\par
        S = 3.267 \quad R-Sq = 96.9\% \quad R-Sq(adj) = 96.7\%\par}

        \vspace{6pt}
        \textbf{(a)} Describe the association between temperature and the number of iced
        coffees sold.\\[4pt]
        \gradedhere{\gradedblank[id=t27a, type=rubric, points=3, rubric=../rubrics/apstats-s01-actual-final-t27a.yaml]{\linewidth+6pt}{\gradedlinesheight{3}}}\ansline{Strong, positive, linear association: days with higher temperatures tend to}\\
        \ansline{have more iced coffees sold. There are no clear outliers.}\\
        \ansline{}\\
        \ansline{}\\
        \ansline{}\\
        \ansline{}\\
        \ansline{}\\
        \ansline{}

        \vspace{6pt}
        \boxguard[16]
        \textbf{(b)} Write the equation of the least-squares regression line, using the
        variable names. Then interpret its slope in context.\\[4pt]
        \gradedhere{\gradedblank[id=t27b, type=rubric, points=3, rubric=../rubrics/apstats-s01-actual-final-t27b.yaml]{\linewidth+6pt}{\gradedlinesheight{3}}}\ansline{predicted coffees $= -51.98 + 1.774$(temperature)}\\
        \ansline{For each additional degree Fahrenheit, the \emph{predicted} number of iced coffees}\\
        \ansline{sold increases by about 1.774 coffees.}\\
        \ansline{}\\
        \ansline{}\\
        \ansline{}\\
        \ansline{}\\
        \ansline{}

        \vspace{6pt}
        \boxguard[21]
        \textbf{(c)} On one day the high temperature was 78$^\circ$F and the bakery sold 80 iced
        coffees. Find the residual for this day and interpret it in context.
        \begin{work}
          \widehat{\text{coffees}} &= -51.98 + 1.774(78) = 86.39 \text{ coffees} \\
          \text{residual} &= y - \hat{y} = 80 - 86.39 = -6.39 \text{ coffees}
        \end{work}
        \strut\gradedhere{\gradedblank[id=t27c, type=rubric, points=3, rubric=../rubrics/apstats-s01-actual-final-t27c.yaml]{\linewidth+6pt}{\gradedlinesheight{3}}}\ansline{On this 78$^\circ$F day the bakery sold about 6.4 \emph{fewer} iced coffees than the}\\
        \ansline{line predicted.}\\
        \ansline{}\\
        \ansline{}\\
        \ansline{}\\
        \ansline{}\\
        \ansline{}\\
        \ansline{}

        \vspace{6pt}
        \boxguard[15]
        \textbf{(d)} Interpret the value of R-Sq in context.\\[4pt]
        \gradedhere{\gradedblank[id=t27d, type=rubric, points=2, rubric=../rubrics/apstats-s01-actual-final-t27d.yaml]{\linewidth+6pt}{\gradedlinesheight{3}}}\ansline{About 96.9\% of the variation in the number of iced coffees sold is explained by}\\
        \ansline{the linear model with daily high temperature.}\\
        \ansline{}\\
        \ansline{}\\
        \ansline{}\\
        \ansline{}\\
        \ansline{}\\
        \ansline{}

        \vspace{6pt}
        \boxguard[17]
        \textbf{(e)} Use the residual plot to decide whether a linear model is appropriate for
        these data. Then explain why the owner should not use this line to predict the number
        of iced coffees sold on a 20$^\circ$F winter day.\\[4pt]
        \gradedhere{\gradedblank[id=t27e, type=rubric, points=4, rubric=../rubrics/apstats-s01-actual-final-t27e.yaml]{\linewidth+6pt}{\gradedlinesheight{3}}}\ansline{Yes: the residuals are scattered randomly around 0 with no curved pattern, so a}\\
        \ansline{linear model is appropriate. But 20$^\circ$F is far below the 62--95$^\circ$F in the}\\
        \ansline{data. That is extrapolation, so the prediction is unreliable: the line gives}\\
        \ansline{about $-16.5$ iced coffees, which is impossible.}\\
        \ansline{}\\
        \ansline{}\\
        \ansline{}\\
        \ansline{}

  % ── FRQ 28 ── Unit 3
  \boxguard[26]
  \item \textbf{(15 points)} A company that makes a language-learning app has 9{,}000
        subscribers. It wants to estimate the proportion of its subscribers who use the app
        every day.

        \vspace{6pt}
        \textbf{(a)} Describe how to use a random number generator to choose a simple random
        sample of 300 subscribers.\\[4pt]
        \gradedhere{\gradedblank[id=t28a, type=rubric, points=3, rubric=../rubrics/apstats-s01-actual-final-t28a.yaml]{\linewidth+6pt}{\gradedlinesheight{3}}}\ansline{Label the subscribers 1 to 9{,}000. Use a random number generator to produce}\\
        \ansline{integers from 1 to 9{,}000, ignoring repeats, until 300 different labels are}\\
        \ansline{chosen. Survey the 300 subscribers with those labels.}\\
        \ansline{}\\
        \ansline{}\\
        \ansline{}\\
        \ansline{}\\
        \ansline{}

        \vspace{6pt}
        \boxguard[17]
        \textbf{(b)} The company sends its survey to the 300 selected subscribers as a
        notification inside the app, and only 120 of them answer. Name the bias this creates.
        Explain who is likely to be under-represented, and whether the estimate is likely too
        high or too low.\\[4pt]
        \gradedhere{\gradedblank[id=t28b, type=rubric, points=3, rubric=../rubrics/apstats-s01-actual-final-t28b.yaml]{\linewidth+6pt}{\gradedlinesheight{3}}}\ansline{Nonresponse bias. Subscribers who rarely open the app are unlikely to see or}\\
        \ansline{answer an in-app survey, so they are under-represented. The estimated proportion}\\
        \ansline{who use the app every day is likely too high.}\\
        \ansline{}\\
        \ansline{}\\
        \ansline{}\\
        \ansline{}\\
        \ansline{}

        \vspace{6pt}
        \boxguard[18]
        \textbf{(c)} The company also wants to know whether a daily reminder notification
        increases practice time. Two hundred subscribers volunteer; 100 will get a daily
        reminder for a month and 100 will get none, and the app will record the minutes each
        one practices that month. Identify the experimental units, the treatments, and the
        response variable.\\[4pt]
        \gradedhere{\gradedblank[id=t28c, type=rubric, points=3, rubric=../rubrics/apstats-s01-actual-final-t28c.yaml]{\linewidth+6pt}{\gradedlinesheight{3}}}\ansline{Experimental units: the 200 volunteer subscribers. Treatments: a daily reminder}\\
        \ansline{notification, or no reminder. Response variable: the minutes each subscriber}\\
        \ansline{practices during the month.}\\
        \ansline{}\\
        \ansline{}\\
        \ansline{}\\
        \ansline{}\\
        \ansline{}

        \vspace{6pt}
        \boxguard[16]
        \textbf{(d)} Describe how to randomly assign the 200 volunteers to the two treatments,
        and explain why random assignment matters in this experiment.\\[4pt]
        \gradedhere{\gradedblank[id=t28d, type=rubric, points=3, rubric=../rubrics/apstats-s01-actual-final-t28d.yaml]{\linewidth+6pt}{\gradedlinesheight{3}}}\ansline{Label the volunteers 1 to 200. Randomly generate 100 different integers from 1 to}\\
        \ansline{200: those volunteers get the reminder, and the other 100 get none. Random}\\
        \ansline{assignment balances other variables (such as motivation) between the groups, so}\\
        \ansline{a difference in practice time can be attributed to the reminders.}\\
        \ansline{}\\
        \ansline{}\\
        \ansline{}\\
        \ansline{}

        \vspace{6pt}
        \boxguard[17]
        \textbf{(e)} The reminder group practiced for significantly more minutes, on average,
        than the no-reminder group. What can the company conclude about cause and effect, and
        to which subscribers does the conclusion apply? Explain.\\[4pt]
        \gradedhere{\gradedblank[id=t28e, type=rubric, points=3, rubric=../rubrics/apstats-s01-actual-final-t28e.yaml]{\linewidth+6pt}{\gradedlinesheight{3}}}\ansline{The treatments were randomly assigned and the difference is significant, so the}\\
        \ansline{company can conclude the reminders caused the extra practice time. But the}\\
        \ansline{subjects were volunteers, not a random sample, so the conclusion applies only to}\\
        \ansline{subscribers like these volunteers, not to all 9{,}000 subscribers.}\\
        \ansline{}\\
        \ansline{}\\
        \ansline{}\\
        \ansline{}

\end{enumerate}

\end{document}
